\documentclass[10pt,a4paper]{amsart}
\usepackage[margin=19mm]{geometry}
\usepackage{amsmath,amssymb,mathtools}
\usepackage[scr=rsfs,cal=euler]{mathalfa}
\usepackage{xfrac}
\usepackage[unicode,hidelinks,bookmarks=false]{hyperref}
\newcommand{\ie}{i.e.\ }
\newcommand{\cf}{cf.\ }
\newcommand{\resp}{respectively\ }
\newcommand{\Subsection}{\subsection}
\providecommand{\nobreakdash}{\nobreak\hbox{-}}
\setlength{\emergencystretch}{2em}
\multlinegap0pt
\usepackage{bm}
\usepackage{bbm}
\usepackage{dsfont}
\usepackage{xspace}
\usepackage{cancel}
\usepackage{mathtools}
\usepackage{amsthm}
\makeatletter
\@ifundefined{theorem}{\newtheorem{theorem}{Theorem}}{}
\@ifundefined{prop}{\newtheorem{prop}[theorem]{Proposition}}{}
\makeatother
\title[Principal Curves In Metric Spaces And The Space Of Probabili]{Principal Curves In Metric Spaces And The Space Of Probability Measures}
\author{Andrew Warren, Anton Afanassiev, Forest Kobayashi, Young-Heon Kim, Geoffrey Schiebinger}
\date{Source: arXiv:2505.04168v2 (CC BY 4.0)}
\begin{document}
\maketitle

\noindent\emph{Source and licence.} Principal Curves In Metric Spaces And The Space Of Probability Measures. Version arXiv:2505.04168v2 (CC BY 4.0), distributed under the Creative Commons Attribution 4.0 International licence, as recorded in the arXiv licence metadata for this exact version and shown on the abstract page https://arxiv.org/abs/2505.04168v2. Full rights notice: licenses/arxiv-2505.04168-CC-BY-4.0.txt.\par\smallskip

\noindent\textit{Editorial scope.} Counted unit: the Prop whose source label is \texttt{prop:discrete-nonlocal-to-continuum} in the article (source0.tex lines 2455--2461, source label \texttt{prop:discrete-nonlocal-to-continuum}), delivered together with the article's own complete original proof of it (source0.tex lines 2493--2609, one \texttt{proof} environment of 117 lines). No mathematics is written, repaired or re-derived: statement and proof are the article's own text line for line. The proof is detached from the statement in the source: the article states the result at lines 2455--2461 and prints its proof at lines 2493--2609, and the proof environment's own header names the statement, so the association is the article's own. Required context retained uncounted: thm:discrete-to-continuum (lines 626--644). Every definition, display and label of those lines is retained unchanged. Omitted from this package and not redistributed: 1--625, 645--2454, 2462--2492, 2610--3192. Nothing inside a retained excerpt is omitted or altered: every retained body is delivered exactly as the article's lines are written, so no omission receipt of any kind accompanies this package. Citations: the retained excerpts carry no citation command at all, so no bibliography entry of the article is transcribed and no reference is redistributed. Reference adaptation: none was needed and none was made; every reference inside the delivered unit resolves inside it, so no unresolved reference or citation is produced. Printed slips kept literal: no printed word, symbol, hypothesis or number of the counted statement or of its proof was altered. Layout and class: the article's own class is \texttt{imsart} with packages including amsthm, amsmath, amsfonts, amssymb, fix-cm, float, silence, mathalfa, natbib, hyperref, graphicx, subcaption, cleveref, crossreftools; the wrapper is the lane's pinned \texttt{amsart} editorial wrapper at 10pt on A4, which restates the theorem-like environments the retained text needs and adds the article's own macro definitions that the retained text needs (none), with extra wrapper packages (bm, bbm, dsfont, xspace, cancel, mathtools). The delivered numbering is the unit's own. Assets: no retained span contains a figure environment, a graphics-inclusion command, a PDF-page inclusion, a table or a TikZ picture; no image file is copied into this package, the package declares no asset dependency and no image file is redistributed with it. Licence: version arXiv:2505.04168v2 of this article is distributed under the Creative Commons Attribution 4.0 International licence, as recorded in the arXiv licence metadata for this exact version and shown on the abstract page https://arxiv.org/abs/2505.04168v2. A full rights notice is delivered at licenses/arxiv-2505.04168-CC-BY-4.0.txt. Characters: every retained excerpt is pure ASCII, so the delivered engine, XeLaTeX with its default Latin Modern text font, reports no missing character.
\begin{theorem}[Discrete to continuum]
\label{thm:discrete-to-continuum}
Let $X$ be a compact geodesic metric space. Let $K,N\in\mathbb{N}$ and
$\Lambda_{N}\in\mathcal{P}(X)$. Suppose also that
$\Lambda_{N}\rightharpoonup^{*}\Lambda$ where
$\Lambda\in\mathcal{P}(X)$. Let $\{\gamma_{k}^{*K}\}_{k=1}^{K}$ be a
minimizer of $\text{PPC}^{K}(\Lambda_{N})$, and let $\gamma^{*K}_t$ be a
constant-speed piecewise geodesic interpolation of
$\{\gamma_{k}^{*K}\}_{k=1}^{K}$.

Then, up to passage to a subsequence
in $K$, we have that: there exists an AC curve $\gamma^{*}$ such that
$\gamma_{t}^{*K}\rightarrow\gamma_{t}^{*}$ uniformly in
$t$, and $\gamma^{*}$ is a minimizer for $\text{PPC}(\Lambda)$.
Moreover, given any uniformly convergent subsequence of geodesic interpolations
$\gamma^{*K}$ of minimizers $\{\gamma_{k}^{*K}\}_{k=1}^{K}$ with limit
$\gamma^{*}\in AC([0,1];X)$, it holds that $\gamma^{*}$ is a minimizer
of $\text{PPC}(\Lambda)$.
\end{theorem}

\begin{prop}
\label{prop:discrete-nonlocal-to-continuum}
Let $X$ be a compact geodesic metric space. Minimizers of the functional
$\text{PPC}_{w}^{K}(\Lambda_{N})$ converge to minimizers of $\text{PPC}(\Lambda)$
as $K,N\rightarrow\infty$ and $\bar{h}\rightarrow0$, in the same
sense as for Theorem \ref{thm:discrete-to-continuum}.
\end{prop}

\begin{proof}[Proof of Proposition \ref{prop:discrete-nonlocal-to-continuum}]
Our proof strategy is to reduce the convergence of minimizers of the nonlocal
discrete functional to the same result for the discrete functional, which was already established in the proof of Theorem \ref{thm:discrete-to-continuum}.

Let  $\{\gamma_1, ..., \gamma_K\}\subset X$.
 For simplicity,  denote for the first term in the functional $\text{PPC}_{w}^{K}(\Lambda_{N})$,
\begin{align*}
  S_{N,K} : =  \frac{1}{N}\sum_{k=1}^{K}\sum_{j=1}^{K}\sum_{x_{n}\in I_{j}}d^{2}(x_{n},\gamma_{k})\frac{1}{C_{j}}w_{h_{j}}\left(\frac{\widehat{\text{Dist}}(\gamma_{j},\gamma_{k})}{\sum_{i=1}^{K-1}d(\gamma_{i},\gamma_{i+1})}\right).
\end{align*}



\emph{Step 1}. From the fact that $I_{j}$
is the Voronoi cell for $\gamma_{j}$, it holds that
\[
\sum_{x_{n}\in I_{j}}d^{2}(x_{n},\gamma_{k})\geq\sum_{x_{n}\in I_{j}}d^{2}(x_{n},\gamma_{j})
\]
and so
\begin{align*}
 S_{N,K} & \geq\frac{1}{N}\sum_{k=1}^{K}\sum_{j=1}^{K}\sum_{x_{n}\in I_{j}}d^{2}(x_{n},\gamma_{j})\frac{1}{C_{j}}w_{h_{j}}\left(\frac{\widehat{\text{Dist}}(\gamma_{j},\gamma_{k})}{\sum_{i=1}^{K-1}d(\gamma_{i},\gamma_{i+1})}\right)\\
 & =\frac{1}{N}\sum_{j=1}^{K}\sum_{x_{n}\in I_{j}}d^{2}(x_{n},\gamma_{j})\frac{1}{C_{j}}\sum_{k=1}^{K}w_{h_{j}}\left(\frac{\widehat{\text{Dist}}(\gamma_{j},\gamma_{k})}{\sum_{i=1}^{K-1}d(\gamma_{i},\gamma_{i+1})}\right).
\end{align*}
Since we chose $C_{j}$ so that
\[
\frac{1}{C_{j}}\sum_{k=1}^{K}w_{h_{j}}\left(\frac{\widehat{\text{Dist}}(\gamma_{j},\gamma_{k})}{\sum_{i=1}^{K-1}d(\gamma_{i},\gamma_{i+1})}\right) \leq 1,
\]
this ensures that for any $\{\gamma_{k}\}_{k=1}^{K}$ and $\{x_{n}\}_{n=1}^{N}$
atoms in $\Lambda_{N}$, and any $\bar{h}$,

\begin{align*}
 S_{N,K}
+\beta\sum_{i=1}^{K-1}d(\gamma_{i},\gamma_{i+1})
\geq\frac{1}{N}\sum_{k=1}^{K}\sum_{x_{n}\in I_{k}}d^{2}(x_{n},\gamma_{k})+\beta\sum_{i=1}^{K-1}d(\gamma_{i},\gamma_{i+1}).
\end{align*}
Accordingly, we can apply Step 1 from the proof of Theorem \ref{thm:discrete-to-continuum},
already established, to deduce that,
 if $\gamma^{K}$, the piecewise geodesic curve from $\{\gamma_1, ..., \gamma_K\}$,  converges uniformly in $t$ to some limiting
AC curve $\gamma$, then
\begin{multline*}
\liminf_{N,K\rightarrow\infty;\bar{h}\rightarrow0}\left(
S_{N,K} +\beta\sum_{i=1}^{K-1}d(\gamma_{i},\gamma_{i+1})\right)\\
\begin{aligned}
& \geq\liminf_{N,K\rightarrow\infty}\left(\frac{1}{N}\sum_{k=1}^{K}\sum_{x_{n}\in I_{k}}d^{2}(x_{n},\gamma_{k})+\beta\sum_{i=1}^{K-1}d(\gamma_{i},\gamma_{i+1})\right)\\
 & \geq\int_{X}d^{2}(x,\Gamma)d\Lambda(x)+ \beta \text{Length}(\gamma).
\end{aligned}
\end{multline*}
(In fact the same inequality holds even if $\bar{h}$ is not sent
to zero; but we do not use this.)


\emph{Step 2}.
We observe that
\[
d^{2}(x_{n},\gamma_{k})\leq d^{2}(x_{n},\gamma_{j})+2d(x_{n},\gamma_{j})d(\gamma_{j},\gamma_{k})+d^{2}(\gamma_{j},\gamma_{k}).
\]
By assumption on $X$, we have that $d(x_{n},\gamma_{j})\leq\text{Diam}(X)$;
and, since we have assumed that  $w_{h_j}$ is compactly supported on $[0,h_j]$,
either
\[
\frac{\widehat{\text{Dist}}(\gamma_{j},\gamma_{k})}{\sum_{i=1}^{K-1}d(\gamma_{i},\gamma_{i+1})}<h_{j},\text{ or }w_{h_{j}}\left(\frac{\widehat{\text{Dist}}(\gamma_{j},\gamma_{k})}{\sum_{i=1}^{K-1}d(\gamma_{i},\gamma_{i+1})}\right)=0.
\]
Additionally, we have that $d(\gamma_{j},\gamma_{k})\leq\widehat{\text{Dist}}(\gamma_{j},\gamma_{k})$
simply from the triangle inequality. Thus, since $h_{j}\leq\bar{h}$,
we have that, for all $j$ and $k$ such that $w_{h_{j}}\left(\frac{\widehat{\text{Dist}}(\gamma_{j},\gamma_{k})}{\sum_{i=1}^{K-1}d(\gamma_{i},\gamma_{i+1})}\right)>0$,
\begin{align*}
d^{2}(x_{n},\gamma_{k}) & \leq d^{2}(x_{n},\gamma_{j})+2d(x_{n},\gamma_{j})\widehat{\text{Dist}}(\gamma_{j},\gamma_{k})+\left(\widehat{\text{Dist}}(\gamma_{j},\gamma_{k})\right)^{2}\\
 & \leq d^{2}(x_{n},\gamma_{j})+2\text{Diam}(X)\bar{h}\sum_{i=1}^{K-1}d(\gamma_{i},\gamma_{i+1})+\left(\bar{h}\sum_{i=1}^{K-1}d(\gamma_{i},\gamma_{i+1})\right)^{2}.
\end{align*}
Now, let us assume that $\sum_{i=1}^{K-1}d(\gamma_{i},\gamma_{i+1})\leq\text{Diam}(X)^{2}/\beta$;
for establishing convergence of minimizers, this shall result in no
loss of generality, as we have already seen in previous arguments.
Under this assumption,
\[
d^{2}(x_{n},\gamma_{k})\leq d^{2}(x_{n},\gamma_{j})+2\text{Diam}(X)^{3}\bar{h}/\beta+\bar{h}^{2}\text{Diam}(X)^{4}/\beta^2.
\]
It follows that, for any $\{\gamma_{k}\}_{k=1}^{K}$, $\Lambda_{N}$,
and $\bar{h}$,

\begin{align*}
 S_{N,K} & \leq\frac{1}{N}\sum_{k=1}^{K}\sum_{j=1}^{K}\sum_{x_{n}\in I_{j}}\left(d^{2}(x_{n},\gamma_{j})+ O(\bar h)
\right)\frac{1}{C_{j}}w_{h_{j}}\left(\frac{\widehat{\text{Dist}}(\gamma_{j},\gamma_{k})}{\sum_{i=1}^{K-1}d(\gamma_{i},\gamma_{i+1})}\right)\\
 & =\frac{1}{N}\sum_{j=1}^{K}\sum_{x_{n}\in I_{j}}\left(d^{2}(x_{n},\gamma_{j})+
 O(\bar h)
 \right)\frac{1}{C_{j}}\sum_{k=1}^{K}w_{h_j}\left(\frac{\widehat{\text{Dist}}(\gamma_{j},\gamma_{k})}{\sum_{i=1}^{K-1}d(\gamma_{i},\gamma_{i+1})}\right)\\
 & \leq \frac{1}{N}\sum_{j=1}^{K}\sum_{x_{n}\in I_{j}}\left(d^{2}(x_{n},\gamma_{j})
 + O(\bar h)
 \right)\\
 & =\left(\frac{1}{N}\sum_{j=1}^{K}\sum_{x_{n}\in I_{j}}d^{2}(x_{n},\gamma_{j})
 \right)+ O(\bar h).
\end{align*}
For the last line, notice that the double summation has only $N$ effective terms.

Consequently,

\begin{multline*}
\limsup_{N,K\rightarrow\infty;\bar{h}\rightarrow0}\left(
 S_{N,K}+\beta\sum_{i=1}^{K-1}d(\gamma_{i},\gamma_{i+1})\right)\\
\begin{aligned}
&
\leq\limsup_{N,K\rightarrow\infty;\bar{h}\rightarrow0}\left(\frac{1}{N}\sum_{j=1}^{K}\sum_{x_{n}\in I_{j}}\left(d^{2}(x_{n},\gamma_{j})+
O(\bar h)
\right)
+\beta\sum_{i=1}^{K-1}d(\gamma_{i},\gamma_{i+1})\right)\\
 & \leq\limsup_{N,K\rightarrow\infty}\left(\frac{1}{N}\sum_{j=1}^{K}\sum_{x_{n}\in I_{j}}d^{2}(x_{n},\gamma_{j})+\beta\sum_{i=1}^{K-1}d(\gamma_{i},\gamma_{i+1})\right).
\end{aligned}
\end{multline*}
Finally, applying Step 2 from the proof of Theorem \ref{thm:discrete-to-continuum}, for the specific choice of $\{\gamma_k\}_{k=1}^K$ from that Step 2, it holds that
\[
\limsup_{N,K\rightarrow \infty}\left(\frac{1}{N}\sum_{j=1}^{K}\sum_{x_{n}\in I_{j}}d^{2}(x_{n},\gamma_{j})+\beta\sum_{i=1}^{K-1}d(\gamma_{i},\gamma_{i+1})\right)\leq\int_{X}d^{2}(x,\Gamma)d\Lambda(x)+\beta\text{Length}(\gamma).
\]

\medskip


 \emph{Step 3.}
Lastly, the conclusion of the proof is now identical to Step 3 of the proof of Theorem \ref{thm:discrete-to-continuum}.
\end{proof}

\end{document}
